%!TEX root = ../graduate-work.tex
\phantomsection
\section*{Введение} 
\addcontentsline{toc}{section}{Введение}

% В этой части надо описать предметную область, задачу из которой вы будете решать, объяснить её актуальность (почему надо что-то делать сейчас?).
% Здесь же стоит ввести определения понятий, которые вам понадобятся в постановке задачи.

% Данная система основывается на $\mathbf{\mathbb {NP}}$-трудной проблеме ... Однако, у неё имеется серьёзный недостаток ...

% В этой работе рассматривается ...  предложенное в 1994 коду~\cite{Ivanov}. В этой работе ...

% В настоящей работе исследуются вопросы ...

% В \S~1 даётся определение ...


% Для некоторых классов приводятся нижние оценки ... (теорема~\ref{theorem1}).

В последние годы наблюдается устойчивый рост интереса к системам автоматической генерации радиологических отчётов на основе изображений медицинских исследований. Современные модели, в том числе трансформерные и мультимодальные архитектуры, демонстрируют всё более высокое качество в задачах генерации клинически осмысленных описаний снимков и рассматриваются как средство снижения нагрузки на специалистов и повышения единообразия отчётности~\cite{pmlr-v158-endo21a,yu2023evaluating,mamdouh2025advancements}. Эти системы позволяют ускорять работу врачей, автоматизировать рутинные части описания и улучшать коммуникацию между специалистами.

Ключевым ограничением для практического внедрения таких систем остаётся надёжная автоматическая оценка качества генерируемых отчётов. Традиционные подходы к оценке текстов, основанные на совпадении N-грамм (последовательностей из N слов) или общем языковом сходстве, плохо отражают клиническую корректность: отчёт, в котором меняется лишь отрицание критически важного признака, может получить высокий балл, хотя с точки зрения врача это серьёзная ошибка. Недавние работы подчёркивают, что подходы к оценке в радиологии должны быть чувствительны к фактическим ошибкам и коррелировать с экспертной разметкой, иначе сравнение моделей и отслеживание прогресса в области генерации радиологических отчётов становится ненадёжным~\cite{yu2023evaluating}.

Для решения этой проблемы появляются специализированные подходы к оценке, использующие клинически информативные представления отчёта. К таким подходам относятся использование векторного представления меток CheXbert для сравнения извлечённых из отчётов патологий~\cite{smit-etal-2020-combining,yu2023evaluating}, а также схему извлечения сущностей и отношений RadGraph, в которой отчёты преобразовываются в граф клинических объектов связей между ними~\cite{jain2021radgraph}. На основе таких представлений можно вычислять показатели совпадения между отчётами, уже непосредственно соотнесённые с клиническими находками, а не с поверхностным лексическим сходством.

Параллельно развиваются методы оценки, опирающиеся на большие языковые модели (large language models, LLM). Подход GREEN использует языковую модель для детального описания и подсчёта клинически значимых и незначимых ошибок в автоматически сгенерированном отчёте-кандидате относительно эталонного отчёта, предлагая количественную оценку и интерпретируемые текстовые объяснения~\cite{ostmeier2024green}. FineRadScore, в свою очередь, формулирует задачу как минимальное число строковых правок, необходимых для превращения отчёта-кандидата в эталонный отчёт, дополняя это оценкой серьёзности ошибок и комментариями к каждой правке~\cite{pmlr-v252-huang24a}. Эти работы показывают, что использование больших языковых моделей позволяет строить более гибкие и интерпретируемые подходы к оценке качества отчётов.

Таким образом, существует острая потребность в разработке специализированного подхода к оценке качества русскоязычных радиологических отчётов, учитывающего особенности языка предметной области и структуру клинических заключений. Целью данной работы является разработка метода оценки качества радиологических отчётов органов грудной клетки на русском языке, который, в отличие от подходов, опирающихся только на текстовое сравнение, учитывает как соответствие отчёта эталонному описанию, так и его согласованность с визуальными данными исследования. Практическая значимость работы определяется возможностью использования предложенного подхода для сопоставления систем автоматической генерации отчётов, контроля качества их результатов и дальнейшей адаптации таких систем к русскоязычной клинической практике.